% Options for packages loaded elsewhere
\PassOptionsToPackage{unicode}{hyperref}
\PassOptionsToPackage{hyphens}{url}
\documentclass[
  11pt,
]{article}
\usepackage{xcolor}
\usepackage[left=3cm,right=5cm,top=2cm,bottom=2cm]{geometry}
\usepackage{amsmath,amssymb}
\setcounter{secnumdepth}{5}
\usepackage{iftex}
\ifPDFTeX
  \usepackage[T1]{fontenc}
  \usepackage[utf8]{inputenc}
  \usepackage{textcomp} % provide euro and other symbols
\else % if luatex or xetex
  \usepackage{unicode-math} % this also loads fontspec
  \defaultfontfeatures{Scale=MatchLowercase}
  \defaultfontfeatures[\rmfamily]{Ligatures=TeX,Scale=1}
\fi
\usepackage{lmodern}
\ifPDFTeX\else
  % xetex/luatex font selection
\fi
% Use upquote if available, for straight quotes in verbatim environments
\IfFileExists{upquote.sty}{\usepackage{upquote}}{}
\IfFileExists{microtype.sty}{% use microtype if available
  \usepackage[]{microtype}
  \UseMicrotypeSet[protrusion]{basicmath} % disable protrusion for tt fonts
}{}
\makeatletter
\@ifundefined{KOMAClassName}{% if non-KOMA class
  \IfFileExists{parskip.sty}{%
    \usepackage{parskip}
  }{% else
    \setlength{\parindent}{0pt}
    \setlength{\parskip}{6pt plus 2pt minus 1pt}}
}{% if KOMA class
  \KOMAoptions{parskip=half}}
\makeatother
\usepackage{longtable,booktabs,array}
\usepackage{caption}
\captionsetup[table]{skip=6pt}
\newcounter{none} % for unnumbered tables
\usepackage{calc} % for calculating minipage widths
% Correct order of tables after \paragraph or \subparagraph
\usepackage{etoolbox}
\makeatletter
\patchcmd\longtable{\par}{\if@noskipsec\mbox{}\fi\par}{}{}
\makeatother
% Allow footnotes in longtable head/foot
\IfFileExists{footnotehyper.sty}{\usepackage{footnotehyper}}{\usepackage{footnote}}
\makesavenoteenv{longtable}
\usepackage{graphicx}
\makeatletter
\newsavebox\pandoc@box
\newcommand*\pandocbounded[1]{% scales image to fit in text height/width
  \sbox\pandoc@box{#1}%
  \Gscale@div\@tempa{\textheight}{\dimexpr\ht\pandoc@box+\dp\pandoc@box\relax}%
  \Gscale@div\@tempb{\linewidth}{\wd\pandoc@box}%
  \ifdim\@tempb\p@<\@tempa\p@\let\@tempa\@tempb\fi% select the smaller of both
  \ifdim\@tempa\p@<\p@\scalebox{\@tempa}{\usebox\pandoc@box}%
  \else\usebox{\pandoc@box}%
  \fi%
}
% Set default figure placement to htbp
\def\fps@figure{htbp}
\makeatother
% definitions for citeproc citations
\NewDocumentCommand\citeproctext{}{}
\NewDocumentCommand\citeproc{mm}{%
  \begingroup\def\citeproctext{#2}\cite{#1}\endgroup}
\makeatletter
 % allow citations to break across lines
 \let\@cite@ofmt\@firstofone
 % avoid brackets around text for \cite:
 \def\@biblabel#1{}
 \def\@cite#1#2{{#1\if@tempswa , #2\fi}}
\makeatother
\newlength{\cslhangindent}
\setlength{\cslhangindent}{1.5em}
\newlength{\csllabelwidth}
\setlength{\csllabelwidth}{3em}
\newenvironment{CSLReferences}[2] % #1 hanging-indent, #2 entry-spacing
 {\begin{list}{}{%
  \setlength{\itemindent}{0pt}
  \setlength{\leftmargin}{0pt}
  \setlength{\parsep}{0pt}
  % turn on hanging indent if param 1 is 1
  \ifodd #1
   \setlength{\leftmargin}{\cslhangindent}
   \setlength{\itemindent}{-1\cslhangindent}
  \fi
  % set entry spacing
  \setlength{\itemsep}{#2\baselineskip}}}
 {\end{list}}
\usepackage{calc}
\newcommand{\CSLBlock}[1]{\hfill\break\parbox[t]{\linewidth}{\strut\ignorespaces#1\strut}}
\newcommand{\CSLLeftMargin}[1]{\parbox[t]{\csllabelwidth}{\strut#1\strut}}
\newcommand{\CSLRightInline}[1]{\parbox[t]{\linewidth - \csllabelwidth}{\strut#1\strut}}
\newcommand{\CSLIndent}[1]{\hspace{\cslhangindent}#1}
\setlength{\emergencystretch}{3em} % prevent overfull lines
\providecommand{\tightlist}{%
  \setlength{\itemsep}{0pt}\setlength{\parskip}{0pt}}
% Three-part paragraph scaffolding for manuscript revision.
%
% bullets  -- boiled-down topic sentence + bullet points
%             small, sans-serif, dark gray; paragraph number [N] prepended
% rgt      -- author's rewrite (initially a placeholder)
%             small, italic, dark blue
% orig     -- original Claude-authored draft
%             normal size, light gray
%
% Presentation order per paragraph: bullets -> rgt -> orig
% Strip with: strip-claudecode -i report.Rmd

\usepackage{xcolor}

\newcounter{pgraph}

\newenvironment{bullets}{%
  \stepcounter{pgraph}%
  \begin{quote}\small\sffamily\color[gray]{0.30}%
  {\normalfont\scriptsize\color[gray]{0.58}[\thepgraph]\enspace}}{%
  \end{quote}}

\newenvironment{rgt}{%
  \begin{quote}\small\itshape\color[RGB]{20,60,160}}{%
  \end{quote}}

\newenvironment{orig}{%
  \begin{quote}\normalsize\color[gray]{0.55}}{%
  \end{quote}}

% backward compat alias
\newenvironment{claudecode}{\begin{orig}}{\end{orig}}
\usepackage{amsmath}
\usepackage{amssymb}
\usepackage{lineno}
\linenumbers
\usepackage{setspace}
\setstretch{1.1}
\usepackage{bookmark}
\IfFileExists{xurl.sty}{\usepackage{xurl}}{} % add URL line breaks if available
\urlstyle{same}
% fallback for those not using the hyperref driver hyperxmp:
\makeatletter
\@ifundefined{xmpquote}{\newcommand{\xmpquote}[1]{#1}}{}
\makeatother
\hypersetup{
  pdftitle={Critical evaluation of the modified Gompertz response function as a parametric template for treatment-response trajectories in N-of-1 trial simulation},
  pdfauthor={pmsimstats team},
  hidelinks,
  pdfcreator={LaTeX via pandoc}}

\title{Critical evaluation of the modified Gompertz response function as a parametric template for treatment-response trajectories in N-of-1 trial simulation}
\author{pmsimstats team}
\date{2026-07-30}

\begin{document}
\maketitle

{
\setcounter{tocdepth}{2}
\tableofcontents
}
\section{Abstract}\label{abstract}

\textbf{Background.} Simulation-based power analyses for aggregated
N-of-1 trials require a parametric template for the within-
participant response trajectory. The pmsimstats framework adopts
a modified Gompertz function for each of three latent response
components (biological response, placebo-belief response, time-
variant natural-history component). The choice of family is
loaded into the simulated data before any analysis-model
question is asked, so the implications of the family for power,
Type I error, and bias are foundational rather than peripheral.

\textbf{Methods.} We evaluate the modified Gompertz family against
the relevant biological, pharmacokinetic, and biostatistical
literature. We characterize the parameter space the function
covers, the trajectories it can and cannot represent, and the
sensitivity of standard pmsimstats inference to misspecification
of the response-curve family. Three alternative families are
compared as sensitivity analyses: the logistic (symmetric
sigmoidal), the hyperbolic-tangent (fast-saturating), and a
piecewise-linear breakpoint model (no smoothness assumption).
For each alternative we compute power, Type I error, and bias of
the biomarker-treatment interaction estimate at trial-relevant
sample sizes against simulated data drawn from the alternative
family but analyzed under the standard linear-mixed model.

\textbf{Results.} A 1\{,\}000-replicate-per-cell simulation across 16
cells (4 trajectory families \(\times\) 2 data-generating process
(DGP) architectures
\(\times\) 2 biomarker effect-size levels) at \(N = 35\), open-label
plus blinded discontinuation (OL+BDC)
design, \(t_{1/2} = 0\), found that the inferential consequence
of trajectory family is architecture-dependent. Under
Architecture B (multivariate normal (MVN) differential
correlation), where the
biomarker-treatment interaction lives in the covariance
structure and is therefore mean-independent, trajectory family
is exactly inert: with matched random-number seeds the four
families return identical rejection decisions, so the
production-seed spread \(\{0.774, 0.777, 0.791, 0.804\}\) is
sampling noise. Under Architecture A (direct mean moderation),
implemented in its faithful multiplicative form in which the
biomarker moderation rides the response trajectory, the four
families separate: powers were \(\{0.743, 0.760, 0.778, 0.782\}\)
for Gompertz, hyperbolic-tangent, piecewise-linear, and logistic
respectively, a 0.039 spread in which the Gompertz default sits
roughly two Monte Carlo standard errors (MCSE) below the logistic and
is the lowest-power, that is mildly conservative, family. Type I
error at the null sat in \([0.019, 0.038]\) across all eight null
cells, slightly conservative and without family-specific
inflation; bias of the interaction estimate was negligible and
family-invariant (\(\approx -0.23\)).

\textbf{Conclusions.} The modified Gompertz family is a defensible
default for the response-curve template in pmsimstats N-of-1
simulation, with two qualifications. Under a covariance-channel
(MVN) data-generating process, trajectory family cannot affect
the interaction test by construction, so the choice is
inferentially free. Under a mean-moderation data-generating
process, trajectory family does carry a modest inferential cost
at the few-percent level, and the Gompertz default is mildly
conservative relative to the symmetric and breakpoint
alternatives. We recommend reporting Gompertz-vs-alternative
sensitivity at least once per parameter regime where simulation
results inform protocol-design decisions, with particular
attention to mean-moderation regimes, and we propose a default
sensitivity grid that the framework can run alongside its
primary cells.

\section{Introduction}\label{introduction}

\subsection{A claim about response-curve choice in N-of-1 trial simulation}\label{a-claim-about-response-curve-choice-in-n-of-1-trial-simulation}

\begin{bullets}

\textbf{The Gompertz function is a defensible default for the
saturating-shape regime in N-of-1 simulation.}

\begin{itemize}\setlength{\itemsep}{2pt}\setlength{\parskip}{0pt}
\item Its asymmetric three-parameter form spans profiles that
  logistic and exponential families cannot represent.
\item Established precedent in tumour dynamics, microbial
  growth, and the unified-Richards family supports the choice.
\item The claim is conditional on the saturating-shape regime
  characteristic of most pharmacodynamic trajectories over
  typical trial timescales.
\end{itemize}

\end{bullets}

\begin{rgt}
Lorem ipsum dolor sit amet, consectetur adipiscing elit, sed do eiusmod tempor incididunt ut labore et dolore magna aliqua.

\end{rgt}

\begin{orig}
The first claim is that the modified Gompertz function used as
the parametric template for treatment-response trajectories in
aggregated N-of-1 trial simulation (specifically, in the
pmsimstats simulation framework that operationalizes the
Hendrickson et al.~design family\textsuperscript{\citeproc{ref-hendrickson2020}{1}}) is a
defensible default for the saturating-shape regime that
characterizes most pharmacodynamic and chronic-condition
response curves over the typical trial timescale. The function
is well-established in tumour growth dynamics\textsuperscript{\citeproc{ref-laird1964dynamics}{2},\citeproc{ref-norton1988gompertzian}{3}}, microbial growth
modeling\textsuperscript{\citeproc{ref-zwietering1990modeling}{4}}, and the broader unified-
Richards growth-curve family\textsuperscript{\citeproc{ref-richards1959flexible}{5}--\citeproc{ref-tjorve2017}{7}}; its three-parameter form is
flexible enough to capture the asymmetric saturation profiles
that linear extrapolation, exponential decay, or symmetric
logistic models cannot.

\end{orig}

\begin{bullets}

\textbf{Family-misspecification sensitivity is absent from the
N-of-1 literature and should be reported routinely.}

\begin{itemize}\setlength{\itemsep}{2pt}\setlength{\parskip}{0pt}
\item Existing N-of-1 simulation studies inherit trajectory
  families without comparing alternatives.
\item Oncology and pharmacology modeling have established
  comparative apparatus that the N-of-1 field has not imported.
\item Routine sensitivity reporting is needed rather than
  silent assumption of family adequacy.
\end{itemize}

\end{bullets}

\begin{rgt}
Lorem ipsum dolor sit amet, consectetur adipiscing elit, sed do eiusmod tempor incididunt ut labore et dolore magna aliqua.

\end{rgt}

\begin{orig}
The second claim is that family-misspecification sensitivity is
essentially absent from the published N-of-1 methodology
literature, and that it should be reported routinely rather
than assumed away. The methodological N-of-1 corpus\textsuperscript{\citeproc{ref-hendrickson2020}{1},\citeproc{ref-lillie2011}{8}--\citeproc{ref-vohra2015cent}{10}} is, to the best of our knowledge, silent on
parametric trajectory shape: papers either inherit a chosen
family from the upstream literature without comparing
alternatives, or specify a linear or AR(1) form that is too
coarse to capture the saturation that real trajectories
exhibit. The oncology-modeling literature has set the
precedent for what a proper family-comparison sensitivity
analysis looks like\textsuperscript{\citeproc{ref-tjorve2010richards}{6},\citeproc{ref-benzekry2014classical}{11}}, and the pharmacology literature has
converged on Hill / sigmoid-Emax / indirect-response models
for related questions\textsuperscript{\citeproc{ref-goutelle2008hill}{12},\citeproc{ref-dayneka1993indirect}{13}}.
The N-of-1 literature has not imported this comparative
apparatus, and the present paper endeavors to supply it.

\end{orig}

\begin{bullets}

\textbf{The case for Gompertz is incomplete without a direct
comparison of inferential consequences across alternative
families.}

\begin{itemize}\setlength{\itemsep}{2pt}\setlength{\parskip}{0pt}
\item The paper does not assert Gompertz is uniquely correct
  but argues it is reasonable within the unified-Richards class.
\item Inferential consequences (power, Type I error, bias)
  under alternative families must be quantified, not asserted.
\item The simulation study uses the standard pmsimstats
  linear-mixed analysis at trial-relevant sample sizes.
\end{itemize}

\end{bullets}

\begin{rgt}
Lorem ipsum dolor sit amet, consectetur adipiscing elit, sed do eiusmod tempor incididunt ut labore et dolore magna aliqua.

\end{rgt}

\begin{orig}
The two claims are coupled because the case for a chosen family
is incomplete without the comparison. We do not argue that
Gompertz is the correct family in any strict sense; rather, we
argue that it is a reasonable choice within the unified-Richards
family of saturating curves, and that the inferential
consequences of choosing it (versus the symmetric logistic, the
von Bertalanffy growth curve\textsuperscript{\citeproc{ref-vonbertalanffy1957}{14}}, or a
piecewise-linear breakpoint model) should be quantified rather
than asserted. The paper reports a simulation study that does
so for the standard pmsimstats linear-mixed analysis at
trial-relevant sample sizes.

\end{orig}

\subsection{Why the question matters}\label{why-the-question-matters}

\begin{bullets}

\textbf{Family choice is a DGP-level decision whose inferential
cost cannot be detected without varying the DGP family while
holding the analysis fixed.}

\begin{itemize}\setlength{\itemsep}{2pt}\setlength{\parskip}{0pt}
\item A single-family simulation cannot reveal the inferential
  cost of that family choice.
\item Robustness to family is a joint property of the DGP and
  analysis specification.
\item This is analysis-model misspecification logic applied one
  level down in the simulation hierarchy.
\end{itemize}

\end{bullets}

\begin{rgt}
Lorem ipsum dolor sit amet, consectetur adipiscing elit, sed do eiusmod tempor incididunt ut labore et dolore magna aliqua.

\end{rgt}

\begin{orig}
First, the chosen family loads into the data before any
analysis-side question is asked. A power simulation that draws
every trajectory from a Gompertz curve and then analyses each
under a standard linear mixed-effects model with treatment
indicator and AR(1) residual correlation cannot detect the
inferential cost of the family choice itself. Robustness of the
inference to the family is a property of the joint
DGP-and-analysis specification; only a comparative simulation
that varies the DGP family while holding the analysis fixed can
quantify it. This is the same logic that underpins
analysis-model misspecification studies in the broader
clinical-trial methodology literature\textsuperscript{\citeproc{ref-morris2019}{15}}, applied
one level down in the simulation hierarchy, and it appears to
us that the N-of-1 simulation literature has not yet drawn the
parallel explicitly.

\end{orig}

\begin{bullets}

\textbf{The Gompertz parameterization has known identifiability
fragility in short noisy trajectories that has not been
characterized for N-of-1 designs.}

\begin{itemize}\setlength{\itemsep}{2pt}\setlength{\parskip}{0pt}
\item Asymptote and rate parameters correlate strongly; the
  shape parameter is near-unidentifiable with few inflection-
  region timepoints.
\item Unified-Richards reparameterizations and sloppiness
  diagnostics exist but have not been applied to N-of-1
  trajectory data.
\item Comparative identifiability results from tumour-growth
  modeling support the need for an analogous N-of-1 diagnosis.
\end{itemize}

\end{bullets}

\begin{rgt}
Lorem ipsum dolor sit amet, consectetur adipiscing elit, sed do eiusmod tempor incididunt ut labore et dolore magna aliqua.

\end{rgt}

\begin{orig}
Second, the function family has known identifiability fragility
in short noisy trajectories. Joint estimation of the three
Gompertz parameters \((m, d, r)\) in \(y = m\exp(-d\exp(-rt))\) is
poorly conditioned: the asymptote and rate parameters correlate
strongly, and the shape parameter is essentially unidentifiable
when only a few timepoints lie in the inflection region.\textsuperscript{\citeproc{ref-tjorve2010richards}{6}} and\textsuperscript{\citeproc{ref-tjorve2017}{7}} propose unified-Richards
reparameterizations in which the parameters become biologically
interpretable (maximum relative growth rate, age at inflection,
value at inflection, asymptote) with lower correlation, and\textsuperscript{\citeproc{ref-jagadeesan2023sloppiness}{16}} develops the formal `sloppiness'
diagnostic that quantifies this directly as a property of the
parameter space.\textsuperscript{\citeproc{ref-benzekry2014classical}{11}} fitted eight
tumour-growth models (exponential, exponential-linear,
power-law, Gompertz, logistic, generalized logistic, von
Bertalanffy, dynamic-carrying-capacity) to two in vivo systems
and ranked them on goodness-of-fit, parametric
identifiability, and predictive power; the power-law model was
the most identifiable in lung data while Gompertz was best for
breast. The same kind of comparative diagnosis has not been
performed for short N-of-1 trajectories, and its absence is,
in our view, problematic.

\end{orig}

\begin{bullets}

\textbf{The Gompertz default in N-of-1 simulation reflects
historical borrowing from growth biology, not empirical
comparison with pharmacology alternatives.}

\begin{itemize}\setlength{\itemsep}{2pt}\setlength{\parskip}{0pt}
\item The pmsimstats Gompertz form traces to tumour dynamics
  and microbial growth literature, not pharmacodynamic evidence.
\item The pharmacology literature defaults to Hill/sigmoid-Emax
  and indirect-response models for comparable questions.
\item The two traditions have not been compared; inheritance
  is not itself evidence of adequacy.
\end{itemize}

\end{bullets}

\begin{rgt}
Lorem ipsum dolor sit amet, consectetur adipiscing elit, sed do eiusmod tempor incididunt ut labore et dolore magna aliqua.

\end{rgt}

\begin{orig}
Third, the default in N-of-1 simulation reflects historical
borrowing rather than empirical optimality. The Gompertz form
entered the pmsimstats framework via the broader growth-curve
tradition, where its use traces to Laird's tumour dynamics work\textsuperscript{\citeproc{ref-laird1964dynamics}{2}} and to Zwietering's microbial growth
reformulations\textsuperscript{\citeproc{ref-zwietering1990modeling}{4}}. The pharmacology
literature, addressing related questions about dose-response
and time-course of pharmacological action, defaults to the
Hill / sigmoid-Emax family\textsuperscript{\citeproc{ref-goutelle2008hill}{12},\citeproc{ref-hill1910}{17}} and
to Jusko's indirect-response models\textsuperscript{\citeproc{ref-dayneka1993indirect}{13}}.
These two traditions have largely not spoken to each other; the
N-of-1 simulation literature has inherited the growth-biology
default without an explicit defense against the pharmacology
default. Either default may be adequate for the saturation
regime, but the inheritance is not itself evidence, and the
question merits a direct empirical answer.

\end{orig}

\subsection{What the published N-of-1 literature does and does not do}\label{what-the-published-n-of-1-literature-does-and-does-not-do}

\begin{bullets}

\textbf{No published N-of-1 simulation study has examined
parametric trajectory-family choice or reported
family-misspecification sensitivity.}

\begin{itemize}\setlength{\itemsep}{2pt}\setlength{\parskip}{0pt}
\item CONSORT-N-of-1, foundational programmatic statements,
  and the Hendrickson framework all omit trajectory-family
  comparison.
\item Silent inheritance of the Gompertz form is pervasive
  across applied N-of-1 simulation work.
\item Adjacent literatures (oncology, pharmacometrics, latent-
  class growth) have comparative apparatus that has not been
  applied to N-of-1 simulation.
\end{itemize}

\end{bullets}

\begin{rgt}
Lorem ipsum dolor sit amet, consectetur adipiscing elit, sed do eiusmod tempor incididunt ut labore et dolore magna aliqua.

\end{rgt}

\begin{orig}
A review of the recent N-of-1 methodology literature confirms,
to the best of our knowledge, a silence on the
parametric-trajectory-family question.

\end{orig}

\begin{bullets}

\textbf{Each major N-of-1 methodological strand silently
inherits the Gompertz default without family comparison.}

\begin{itemize}\setlength{\itemsep}{2pt}\setlength{\parskip}{0pt}
\item CONSORT-N-of-1 reporting standards, Lillie and Schork
  programmatic statements, and the Hendrickson framework all
  omit family comparison.
\item Applied prazosin-PTSD work inherits the framework without
  revisiting the trajectory assumption.
\item No published N-of-1 simulation study, to the authors'
  knowledge, has broached the misspecification question.
\end{itemize}

\end{bullets}

\begin{rgt}
Lorem ipsum dolor sit amet, consectetur adipiscing elit, sed do eiusmod tempor incididunt ut labore et dolore magna aliqua.

\end{rgt}

\begin{orig}
The CONSORT extension for N-of-1 trials\textsuperscript{\citeproc{ref-vohra2015cent}{10}}
codifies reporting standards but does not address parametric
trajectory shape. The foundational programmatic statements\textsuperscript{\citeproc{ref-lillie2011}{8},\citeproc{ref-schork2023precision}{9}} emphasize design and
analysis broadly without committing to a parametric response
form. The Hendrickson et al.~framework\textsuperscript{\citeproc{ref-hendrickson2020}{1}}
introduces the modified Gompertz template into the simulation
engine for biomarker-validation power analyses without
comparing it against alternatives. Recent applied papers in
chronic-condition N-of-1 trials (including the prazosin-PTSD
reference work building on\textsuperscript{\citeproc{ref-raskind2013prazosin}{18}}) inherit the
framework without examining the choice. Unfortunately, this
pattern of silent inheritance is pervasive enough that no
published N-of-1 simulation study, to the best of our
knowledge, has broached the family-misspecification question.

\end{orig}

\begin{bullets}

\textbf{Comparative trajectory-family apparatus exists in
adjacent literatures but has not been applied to N-of-1
simulation.}

\begin{itemize}\setlength{\itemsep}{2pt}\setlength{\parskip}{0pt}
\item Oncology modeling, pharmacometrics, and latent-class
  growth literature each have established multi-family
  comparison methods.
\item None of these methods has been applied with the Gompertz
  form as the N-of-1 simulation reference.
\item Importing this machinery is the explicit contribution of
  the present paper.
\end{itemize}

\end{bullets}

\begin{rgt}
Lorem ipsum dolor sit amet, consectetur adipiscing elit, sed do eiusmod tempor incididunt ut labore et dolore magna aliqua.

\end{rgt}

\begin{orig}
In adjacent literatures the comparative apparatus does exist.
The oncology-modeling community routinely fits multiple growth
families and compares them on identifiability and predictive
performance\textsuperscript{\citeproc{ref-benzekry2014classical}{11}}. The pharmacometrics
community produces explicit sensitivity studies for sigmoid-Emax
and indirect-response models\textsuperscript{\citeproc{ref-ebling1996feasibility}{19},\citeproc{ref-triggs1996levodopa}{20}}. The latent-class growth-trajectory
literature treats response-curve heterogeneity through mixture
modeling\textsuperscript{\citeproc{ref-hockenberry2017symptom}{21}}. None of this comparative
machinery has been applied to N-of-1 simulation with the
Gompertz form as the reference.

\end{orig}

\begin{bullets}

\textbf{The paper's contribution is a four-family comparison
against the standard pmsimstats analysis, reported under the
ADEMP framework.}

\begin{itemize}\setlength{\itemsep}{2pt}\setlength{\parskip}{0pt}
\item Four families are compared: Gompertz, symmetric logistic,
  hyperbolic-tangent, and piecewise-linear breakpoint.
\item The analysis model is held fixed at the standard
  pmsimstats linear-mixed specification.
\item Results are reported in the Morris-White-Crowther ADEMP
  framework adopted as the project-wide standard.
\end{itemize}

\end{bullets}

\begin{rgt}
Lorem ipsum dolor sit amet, consectetur adipiscing elit, sed do eiusmod tempor incididunt ut labore et dolore magna aliqua.

\end{rgt}

\begin{orig}
The present paper's contribution is to import that machinery,
specifically: a four-family comparison (Gompertz, symmetric
logistic, hyperbolic-tangent, piecewise-linear breakpoint)
crossed against the standard pmsimstats linear-mixed analysis,
with results reported in the Morris-White-Crowther\textsuperscript{\citeproc{ref-morris2019}{15}} ADEMP framework already adopted as the
project-wide simulation-reporting standard.

\end{orig}

\subsection{What this paper contributes}\label{what-this-paper-contributes}

\begin{bullets}

\textbf{Section 2 characterizes the pmsimstats Gompertz
parameterization and its relationship to the unified-Richards
family.}

\begin{itemize}\setlength{\itemsep}{2pt}\setlength{\parskip}{0pt}
\item The vertical-offset adjustment and calibration parameter
  sets are documented.
\item The unified-Richards reparameterization provides more
  interpretable, less correlated parameters.
\item This section establishes the reference point for all
  subsequent sensitivity comparisons.
\end{itemize}

\end{bullets}

\begin{rgt}
Lorem ipsum dolor sit amet, consectetur adipiscing elit, sed do eiusmod tempor incididunt ut labore et dolore magna aliqua.

\end{rgt}

\begin{orig}
In section 2 we characterize the modified Gompertz family as
implemented in pmsimstats: the vertical-offset adjustment, the
parameter regions corresponding to the calibration sets used
across published pmsimstats simulation studies, and the
unified-Richards reparameterization\textsuperscript{\citeproc{ref-tjorve2010richards}{6},\citeproc{ref-tjorve2017}{7}} that supplies more interpretable and less
correlated parameters.

\end{orig}

\begin{bullets}

\textbf{Section 3 surveys three alternative families at matched
marginal effect size for visual and structural comparison with
Gompertz.}

\begin{itemize}\setlength{\itemsep}{2pt}\setlength{\parskip}{0pt}
\item Families covered: symmetric logistic, hyperbolic-tangent,
  piecewise-linear breakpoint.
\item Each is calibrated to match the Gompertz reference at the
  same asymptote and anchor-point response level.
\item Shape differences (symmetry, curvature, smoothness) are
  the primary axes of contrast.
\end{itemize}

\end{bullets}

\begin{rgt}
Lorem ipsum dolor sit amet, consectetur adipiscing elit, sed do eiusmod tempor incididunt ut labore et dolore magna aliqua.

\end{rgt}

\begin{orig}
In section 3 we survey three alternative response-curve
families: the symmetric logistic\textsuperscript{\citeproc{ref-verhulst1838}{22}}, the
hyperbolic-tangent (formally a rescaled symmetric logistic), and
a piecewise-linear breakpoint model. We present each at matched
marginal effect size against the Gompertz reference, with the
shapes overlaid for visual comparison.

\end{orig}

\begin{bullets}

\textbf{Sections 4 and 5 report a Monte Carlo sensitivity
simulation comparing power, Type I error, and bias across the
four trajectory families.}

\begin{itemize}\setlength{\itemsep}{2pt}\setlength{\parskip}{0pt}
\item Each family generates trajectories; all are analyzed
  under the fixed standard pmsimstats linear-mixed model.
\item Performance measures are reported with conventional MCSEs.
\item The design borrows the Benzekry eight-model comparative
  methodology and the Jagadeesan sloppiness diagnostic.
\end{itemize}

\end{bullets}

\begin{rgt}
Lorem ipsum dolor sit amet, consectetur adipiscing elit, sed do eiusmod tempor incididunt ut labore et dolore magna aliqua.

\end{rgt}

\begin{orig}
In sections 4 and 5 we conduct a Monte Carlo sensitivity-analysis
simulation that generates trajectories under each of the four
families and analyses them under the standard pmsimstats
linear-mixed model. We report power, Type I error, and bias of
the biomarker-treatment interaction estimate under each cell,
with conventional MCSEs. The simulation
borrows the comparative methodology of\textsuperscript{\citeproc{ref-benzekry2014classical}{11}}
(eight-model tumour-growth comparison) and the sloppiness
diagnostic of\textsuperscript{\citeproc{ref-jagadeesan2023sloppiness}{16}} for small-sample
identifiability assessment.

\end{orig}

\begin{bullets}

\textbf{Section 6 identifies the trajectory regimes where
Gompertz misspecification would be inferentially consequential
and recommends sensitivity-grid extensions.}

\begin{itemize}\setlength{\itemsep}{2pt}\setlength{\parskip}{0pt}
\item Consequential regimes include non-saturating growth,
  biphasic patterns, and breakpoint behavior.
\item Standard N-of-1 inference is materially affected by
  family choice in these regimes.
\item Specific sensitivity-grid extensions are recommended for
  pmsimstats users.
\end{itemize}

\end{bullets}

\begin{rgt}
Lorem ipsum dolor sit amet, consectetur adipiscing elit, sed do eiusmod tempor incididunt ut labore et dolore magna aliqua.

\end{rgt}

\begin{orig}
In section 6 we identify the trajectory regimes in which
Gompertz misspecification would be expected to be consequential:
non-saturating growth beyond the trial window, biphasic patterns,
and breakpoint behavior. For these regimes, the standard N-of-1
inference is materially affected by the family choice, and we
recommend specific sensitivity-grid extensions that pmsimstats
users should run alongside their primary cells.

\end{orig}

\begin{bullets}

\textbf{A companion technical document preserves full
mathematical and historical context omitted from the manuscript
proper.}

\begin{itemize}\setlength{\itemsep}{2pt}\setlength{\parskip}{0pt}
\item Location: \texttt{docs/25-gompertz-model-evaluation.md}.
\item Contains citations, derivations, and historical context
  that did not survive manuscript editing.
\item The manuscript is self-contained; the companion document
  is supplementary reference material.
\end{itemize}

\end{bullets}

\begin{rgt}
Lorem ipsum dolor sit amet, consectetur adipiscing elit, sed do eiusmod tempor incididunt ut labore et dolore magna aliqua.

\end{rgt}

\begin{orig}
The full mathematical and historical context, including citations
that did not survive editing into the manuscript proper, is
preserved in the companion technical document at
\texttt{docs/25-gompertz-model-evaluation.md}.

\end{orig}

\section{The modified Gompertz function in pmsimstats}\label{the-modified-gompertz-function-in-pmsimstats}

\begin{bullets}

\textbf{The pmsimstats framework adopts a modified monotone-
saturating Gompertz function for each of three latent response
components.}

\begin{itemize}\setlength{\itemsep}{2pt}\setlength{\parskip}{0pt}
\item Provenance traces to Gompertz's 1825 mortality law and
  Winsor's 1932 growth-curve application.
\item Three components are parameterized: biological response
  BR, placebo-belief response PB, and natural-history TV.
\item The modified form is monotone-saturating, distinguishing
  it from the original mortality-curve application.
\end{itemize}

\end{bullets}

\begin{rgt}
Lorem ipsum dolor sit amet, consectetur adipiscing elit, sed do eiusmod tempor incididunt ut labore et dolore magna aliqua.

\end{rgt}

\begin{orig}
The Gompertz curve originates in Gompertz's 1825 law of human
mortality\textsuperscript{\citeproc{ref-gompertz1825}{23}} and was first applied as a general
growth curve by\textsuperscript{\citeproc{ref-winsor1932gompertz}{24}}; the pmsimstats framework
adopts the modified, monotone-saturating form of that curve.
The framework parameterizes each of the three latent response
components (biological response BR, placebo-belief response PB,
and time-variant natural-history component TV) with a modified
Gompertz function

\end{orig}

\begin{equation}
y(t) = \mathrm{maxr} \cdot \exp\!\left(-\mathrm{disp} \cdot
       \exp(-\mathrm{rate} \cdot t)\right),
\label{eq:gompertz}
\end{equation}

\begin{bullets}

\textbf{The three-parameter Gompertz form is asymmetric and
monotone-saturating, with the inflection location determined
jointly by rate and displacement.}

\begin{itemize}\setlength{\itemsep}{2pt}\setlength{\parskip}{0pt}
\item Parameters are maxr (asymptote), rate (approach speed),
  and disp (curvature near $t = 0$); a vertical offset ensures
  $y(0) = 0$.
\item Inflection at $t^* = \log(\mathrm{disp})/\mathrm{rate}$;
  the asymmetry is the key structural distinction from logistic.
\item Prazosin/PTSD reference values are
  $\mathrm{maxr}_{BR} \approx 4$, $\mathrm{rate}_{BR}
  \approx 0.5$, $\mathrm{disp}_{BR} \approx 4$.
\end{itemize}

\end{bullets}

\begin{rgt}
Lorem ipsum dolor sit amet, consectetur adipiscing elit, sed do eiusmod tempor incididunt ut labore et dolore magna aliqua.

\end{rgt}

\begin{orig}
with vertical-offset adjustment so that \(y(0) = 0\). The three
parameters are \(\mathrm{maxr}\) (asymptotic maximum response,
matching the on-drug or on-natural-history saturation level),
\(\mathrm{rate}\) (which controls how quickly the response
approaches its asymptote), and \(\mathrm{disp}\) (which controls
the curvature near \(t = 0\)). The Gompertz form is
monotone-saturating and asymmetric: the inflection occurs at
\(t^* = \log(\mathrm{disp}) / \mathrm{rate}\), with the rise
being relatively slower before the inflection and faster after.
The parameterization used at the prazosin/PTSD reference cell
sets \(\mathrm{maxr}_{BR} \approx 4\) nightmares per week,
\(\mathrm{rate}_{BR} \approx 0.5\) per week, and
\(\mathrm{disp}_{BR} \approx 4\) for biological response;
analogous parameter sets for PB and TV are documented in
\texttt{docs/25-gompertz-model-evaluation.md} and
\texttt{R/utilities.R}.

\end{orig}

\section{Alternative response-curve families}\label{alternative-response-curve-families}

\begin{bullets}

\textbf{Three alternative families are calibrated to match
Gompertz at a common anchor point, isolating shape effects
from amplitude effects.}

\begin{itemize}\setlength{\itemsep}{2pt}\setlength{\parskip}{0pt}
\item Families compared: symmetric logistic, hyperbolic-tangent,
  piecewise-linear breakpoint.
\item Calibration constraint: matched asymptote and response
  level at $t = 5$ weeks via uniroot or closed form.
\item Families differ in smoothness, inflection symmetry, and
  curvature while sharing the same marginal effect size.
\end{itemize}

\end{bullets}

\begin{rgt}
Lorem ipsum dolor sit amet, consectetur adipiscing elit, sed do eiusmod tempor incididunt ut labore et dolore magna aliqua.

\end{rgt}

\begin{orig}
Three alternative families serve as sensitivity-analysis
comparators, and in this section we shall describe each in turn
before the simulation study proper. Each is parameterized so
that the asymptotic maximum and a single-anchor calibration
point (the response level at \(t = 5\) weeks, calibrated to the
Gompertz family at the same anchor) are matched across families.
That is, the families differ in shape envelope (smoothness,
symmetry around the inflection, and curvature) while sharing
the same marginal effect size at the anchor, a design choice
that we discuss further in section 4.

\end{orig}

\subsection{Symmetric logistic}\label{symmetric-logistic}

The logistic function

\begin{equation}
y(t) = \frac{\mathrm{maxr}}{1 + \exp(-\mathrm{rate} \cdot
              (t - t_0))} - \mathrm{offset}
\label{eq:logistic}
\end{equation}

\begin{bullets}

\textbf{The symmetric logistic rises faster than Gompertz
before the inflection and slower after, with symmetric
inflection as the principal structural contrast.}

\begin{itemize}\setlength{\itemsep}{2pt}\setlength{\parskip}{0pt}
\item Symmetric around inflection at $t_0$; a vertical offset
  ensures $y(0) = 0$.
\item At matched asymptote and anchor, the logistic envelope
  differs visibly from Gompertz in pre- and post-inflection
  curvature.
\item Inflection symmetry is the primary qualitative contrast
  with the asymmetric Gompertz.
\end{itemize}

\end{bullets}

\begin{rgt}
Lorem ipsum dolor sit amet, consectetur adipiscing elit, sed do eiusmod tempor incididunt ut labore et dolore magna aliqua.

\end{rgt}

\begin{orig}
is monotone-saturating and symmetric around its inflection at
\(t = t_0\). A vertical offset is subtracted so that \(y(0) = 0\).
At matched \(\mathrm{maxr}\) and matched response level at the
calibration anchor, the logistic family rises faster than
Gompertz before the inflection and slower after, producing a
visibly different shape envelope while sharing the same
asymptote. The symmetric inflection is the principal qualitative
contrast with the asymmetric Gompertz curve, and the reader will
note that this symmetry property is, in a sense, the logistic
family's most theoretically interesting feature for the purposes
of the present comparison.

\end{orig}

\subsection{Hyperbolic-tangent}\label{hyperbolic-tangent}

The hyperbolic-tangent variant

\begin{equation}
y(t) = \mathrm{maxr} \cdot \frac{\tanh(\mathrm{rate} \cdot
       (t - t_0)) + 1}{2} - \mathrm{offset}
\label{eq:tanh}
\end{equation}

\begin{bullets}

\textbf{The hyperbolic-tangent family represents the fast-
saturating smooth regime that neither the logistic nor Gompertz
can span with standard parameterizations.}

\begin{itemize}\setlength{\itemsep}{2pt}\setlength{\parskip}{0pt}
\item Symmetric and monotone-saturating, but with a steeper
  pre-asymptote rise than the logistic.
\item Used in pharmacokinetic modeling for drugs with rapid
  receptor saturation.
\item Its inclusion spans the fast-saturating smooth regime
  absent from the other three families.
\end{itemize}

\end{bullets}

\begin{rgt}
Lorem ipsum dolor sit amet, consectetur adipiscing elit, sed do eiusmod tempor incididunt ut labore et dolore magna aliqua.

\end{rgt}

\begin{orig}
is also monotone-saturating and symmetric, but with a
qualitatively faster approach to the asymptote than the
logistic. The tanh family produces the steepest pre-asymptote
rise of the three smooth families, and it has been used in
pharmacokinetic modeling for drugs with rapid receptor
saturation. We include it here as the representative of the
fast-saturating smooth regime that neither the logistic nor the
Gompertz family can span with their standard parameterizations.

\end{orig}

\subsection{Piecewise-linear breakpoint}\label{piecewise-linear-breakpoint}

The piecewise-linear-breakpoint family, formalized in the
broken-line regression literature\textsuperscript{\citeproc{ref-muggeo2003}{25}},

\begin{equation}
y(t) = \mathrm{maxr} \cdot \min(1,\, t / t_{\mathrm{bp}})
       + \mathrm{post\_slope} \cdot \max(0,\, t - t_{\mathrm{bp}})
\label{eq:plb}
\end{equation}

\begin{bullets}

\textbf{The piecewise-linear breakpoint is the canonical non-
smooth alternative and provides the strongest test of the
framework's sensitivity to trajectory non-smoothness.}

\begin{itemize}\setlength{\itemsep}{2pt}\setlength{\parskip}{0pt}
\item A linear ramp to breakpoint $t_\mathrm{bp}$ then
  constant or small post-breakpoint slope.
\item The breakpoint is the sole non-differentiable feature
  across the four families.
\item Whether the Wald test detects the non-smooth feature is
  the central empirical question the simulation addresses.
\end{itemize}

\end{bullets}

\begin{rgt}
Lorem ipsum dolor sit amet, consectetur adipiscing elit, sed do eiusmod tempor incididunt ut labore et dolore magna aliqua.

\end{rgt}

\begin{orig}
is the canonical non-smooth alternative: a linear ramp up to a
breakpoint \(t_{\mathrm{bp}}\), then constant (when
\(\mathrm{post\_slope} = 0\)) or with a small post-breakpoint
slope. The breakpoint is the only non-differentiable feature
across the four families and is the principal qualitative
contrast with the three smooth families. This family is the
strongest available test of the framework's sensitivity to
non-smoothness in the trajectory. But how different is the
piecewise-linear breakpoint from the smooth families, from the
perspective of the Wald test that the analysis model applies?
That question is precisely what the simulation in section 4
is designed to answer.

\end{orig}

\section{Sensitivity-analysis simulation design}\label{sensitivity-analysis-simulation-design}

\begin{bullets}

\textbf{The primary estimand is the bm:Dbc interaction
coefficient; performance measures are power, Type I error,
bias, and empirical MCSE.}

\begin{itemize}\setlength{\itemsep}{2pt}\setlength{\parskip}{0pt}
\item Analysis model: $\mathrm{Sx} \sim bm + t + Dbc +
  bm\!:\!Dbc + (1 \mid \mathrm{ptID})$ with
  \texttt{corCAR1(form = \textasciitilde{}t|ptID)}.
\item Primary outputs: power at $\alpha = 0.05$ and Type I
  error under $\beta_{bm:D} = 0$.
\item Secondary outputs: bias and empirical MCSE of
  $\hat\beta_{bm:D}$; all proportions reported with
  binomial MCSE.
\end{itemize}

\end{bullets}

\begin{rgt}
Lorem ipsum dolor sit amet, consectetur adipiscing elit, sed do eiusmod tempor incididunt ut labore et dolore magna aliqua.

\end{rgt}

\begin{orig}
The primary estimand is the biomarker-treatment interaction
coefficient \(\hat{\beta}_{bm:D}\) from the linear mixed-effects
analysis model
\(\mathrm{Sx} \sim bm + t + Dbc + bm\!:\!Dbc + (1 \mid
\mathrm{ptID})\)
with \texttt{corCAR1(form = \textasciitilde{}t|ptID)}.
Primary inferential outputs are
power for \(H_0: \beta_{bm:D} = 0\) at \(\alpha = 0.05\), and Type
I error under \(\beta_{bm:D} = 0\). Secondary estimands are bias
of \(\hat{\beta}_{bm:D}\) relative to the calibrated true effect
and the empirical MCSE of
\(\hat{\beta}_{bm:D}\) across replicates. Performance measures are
reported with binomial-proportion MCSEs for
proportions and the standard MCSE of the
within-replicate mean for continuous quantities.

\end{orig}

\begin{bullets}

\textbf{The ADEMP design is pre-registered and fixes all four
DGP families, the calibration sub-study, and the full factorial
before production runs.}

\begin{itemize}\setlength{\itemsep}{2pt}\setlength{\parskip}{0pt}
\item Pre-registration location:
  \texttt{analysis/scripts/gompertz-evaluation/00-ademp-pre-registration.md}.
\item Study 0 calibrates all four families to a common marginal
  effect size; Study 1 is the 144-cell factorial; Study 2 is
  the sloppiness eigenvalue diagnostic.
\item Pre-registration prevents post-hoc specification of
  performance measures or DGP choices.
\end{itemize}

\end{bullets}

\begin{rgt}
Lorem ipsum dolor sit amet, consectetur adipiscing elit, sed do eiusmod tempor incididunt ut labore et dolore magna aliqua.

\end{rgt}

\begin{orig}
The full Aims, Data-generating mechanisms, Estimands, Methods,
and Performance measures are pre-registered at
\texttt{analysis/scripts/gompertz-evaluation/00-ademp-pre-registration.md},
which fixes the four DGP families (modified Gompertz, symmetric
logistic, hyperbolic-tangent, piecewise-linear breakpoint), the
calibration sub-study (Study 0) that matches all four to a
common marginal effect size, the four-family power-and-bias
factorial (Study 1, 144 cells), and the sloppiness eigenvalue
diagnostic (Study 2) before the production runs.

\end{orig}

\begin{bullets}

\textbf{The present submission uses a focused 16-cell factorial
scoped to answer whether trajectory family matters for the
bm:Dbc Wald test at MCSE 0.013--0.016.}

\begin{itemize}\setlength{\itemsep}{2pt}\setlength{\parskip}{0pt}
\item Design: 4 families $\times$ 2 architectures $\times$
  2 effect sizes; prazosin-calibrated OL+BDC, $N = 35$,
  $t_{1/2} = 0$.
\item Carryover is absent to isolate trajectory shape from
  carryover-mediated correlation decay.
\item The 144-cell extended grid is deferred to a follow-up;
  the 16-cell factorial addresses the central question at
  adequate MCSE resolution.
\end{itemize}

\end{bullets}

\begin{rgt}
Lorem ipsum dolor sit amet, consectetur adipiscing elit, sed do eiusmod tempor incididunt ut labore et dolore magna aliqua.

\end{rgt}

\begin{orig}
For the present submission the simulation is scoped to a 16-cell
focused factorial: four trajectory families crossed with two DGP
architectures (mean moderation and MVN differential correlation,
following the companion manuscript at
\texttt{analysis/report/01-dgp-mean-moderation-vs-mvn/})
crossed with two biomarker effect-size levels
(\(c_{bm} \in \{0, 0.45\}\)). The trial design is the
prazosin-calibrated Hybrid OL+BDC, \(N = 35\), \(t_{1/2} = 0\),
with carryover set to absent in order to isolate the
trajectory-shape effect from carryover-mediated correlation
decay. Each cell is run at \(n_{\text{reps}} = 1{,}000\)
replicates with 8-core \texttt{parallel::mclapply}. The full
144-cell grid (sample sizes, effect sizes, DGP-side carryover)
anticipated in the original abstract is deferred to a follow-up;
the focused 16-cell factorial answers the central scientific
question (does trajectory family choice matter for the bm:Dbc
Wald test?) at MCSE 0.013--0.016 per cell.

\textbf{Notation for the moderation parameter.} Because this factorial
crosses both architectures, \(c_{bm}\) is used here as a common label
for the biomarker-moderation strength at a matched numeric value.
Strictly, the parameter is the correlation \(c_{bm}\) under
Architecture B and the dimensionless multiplier \(\beta_{bm}\) under
Architecture A, the two being calibrated to a common scale because
the Architecture-A shift is applied as \(\beta_{bm}\,b_i\,\sigma_{BR}\)
on the standardized biomarker \(b_i\). The companion architecture
manuscript keeps the symbols distinct; we adopt the shared label
only where a statement applies to both architectures at the matched
value, and the architecture is named explicitly wherever a result
is architecture-specific.

\end{orig}

\begin{bullets}

\textbf{All four families are matched at $t = 5$ weeks with
the asymptotic maximum held constant; calibrated parameters
are recorded for reproducibility.}

\begin{itemize}\setlength{\itemsep}{2pt}\setlength{\parskip}{0pt}
\item Anchor calibration: each family's response at $t = 5$
  matched to Gompertz via uniroot or closed form.
\item Asymptotic maximum held constant across all four families.
\item Calibrated parameters recorded in
  \texttt{02-faithful-trajectory-sweep.R} and transcribed into
  the summary at \texttt{07-gompertz-faithful-summary.txt}.
\end{itemize}

\end{bullets}

\begin{rgt}
Lorem ipsum dolor sit amet, consectetur adipiscing elit, sed do eiusmod tempor incididunt ut labore et dolore magna aliqua.

\end{rgt}

\begin{orig}
Calibration. The four trajectory families are matched at a
single anchor point (\(t = 5\) weeks): each family's response
level at \(t = 5\) is set equal to the Gompertz reference level
at \(t = 5\) via uniroot (logistic, hyperbolic-tangent) or closed
form (piecewise-linear breakpoint). The asymptotic maximum is
held constant across families. Calibrated parameters are
recorded in the driver
\texttt{analysis/scripts/gompertz-evaluation/02-faithful-trajectory-sweep.R}.

\end{orig}

\begin{bullets}

\textbf{The trajectory family drives the data-generating process
directly, and under Architecture A the biomarker moderation is
trajectory-scaled, in keeping with the paper-01 multiplicative
model.}

\begin{itemize}\setlength{\itemsep}{2pt}\setlength{\parskip}{0pt}
\item The family sets the BR mean curve inside the covariance
  construction (\texttt{buildSigma} \texttt{br\_family}), not by
  a post-hoc additive shift.
\item Under Architecture A the per-timepoint moderation scales
  with the response trajectory, normalised so the mean on-drug
  magnitude equals $\sigma_{BR}$ (matched marginal effect).
\item Under Architecture B the interaction is in the covariance
  and is mean-independent, so the family enters the BR mean only.
\end{itemize}

\end{bullets}

\begin{rgt}
Lorem ipsum dolor sit amet, consectetur adipiscing elit, sed do eiusmod tempor incididunt ut labore et dolore magna aliqua.

\end{rgt}

\begin{orig}
Trajectory-family encoding. The family is routed through the
data-generating process rather than imposed post hoc: it sets
the BR mean curve inside the covariance construction via the
\texttt{br\_family} argument to \texttt{buildSigma}. Under
Architecture A (direct mean moderation) the biomarker-treatment
interaction is implemented in the multiplicative form of the
companion manuscript at
\texttt{analysis/report/01-dgp-mean-moderation-vs-mvn/},
\(Y_{it} = \mu_0 + \beta_1 D_{it}(1 + \beta_{bm} B_i) +
\epsilon_{it}\), so the biomarker-dependent component of the
treatment effect rides the response trajectory \(BR(t)\) rather
than entering as a flat shift. We normalise the per-timepoint
moderation so that its mean on-drug magnitude equals
\(\sigma_{BR}\), holding the marginal interaction effect size
fixed across families while letting the temporal profile vary by
family shape. Under Architecture B (MVN differential
correlation) the interaction is encoded in the BM-BR covariance
block, which is independent of the BR mean; the family therefore
enters only through the mean curve and, as we verify by a
matched-seed check, cannot move the interaction test.

\end{orig}

\section{Results}\label{results}

\begin{bullets}

\textbf{The 16-cell simulation completed with 100\% convergence
in 597 s; cell MCSE on power runs from 0.013 to 0.014.}

\begin{itemize}\setlength{\itemsep}{2pt}\setlength{\parskip}{0pt}
\item 1{,}000 replicates per cell across 16 cells;
  8-core \texttt{parallel::mclapply}.
\item Per-replicate output at
  \texttt{analysis/data/quick-sim/07-gompertz-faithful-replicates.rds};
  ADEMP summary at \texttt{07-gompertz-faithful-summary.txt}.
\item Results are presented in order: power, Type I error,
  bias.
\end{itemize}

\end{bullets}

\begin{rgt}
Lorem ipsum dolor sit amet, consectetur adipiscing elit, sed do eiusmod tempor incididunt ut labore et dolore magna aliqua.

\end{rgt}

\begin{orig}
The simulation program was executed at 1\{,\}000 replicates per
cell across the 16-cell trajectory-family-by-architecture-by-
effect-size factorial in 597 s using 8-core
\texttt{parallel::mclapply}.
Convergence was 100\% across all 16\{,\}000 fits. Per-replicate
output is at
\texttt{analysis/data/quick-sim/07-gompertz-faithful-replicates.rds};
the Morris ADEMP cell-level summary is at
\texttt{analysis/data/quick-sim/07-gompertz-faithful-summary.txt}.
Cell MCSE on power runs from 0.013 to 0.014 across the
alternative cells and from 0.004 to 0.006 on the conservative
null cells. We begin the results section with power, then turn
to Type I error, and finally consider bias of the interaction
estimate.

\end{orig}

\subsection{Power across response-curve families}\label{power-across-response-curve-families}

Trajectory-family power at \(c_{bm} = 0.45\), percent \(\pm\) MCSE.
Architecture A is the faithful trajectory-scaled mean-moderation
form; Architecture B is the MVN covariance channel:

{\def\LTcaptype{none} % do not increment counter
\begin{longtable}[]{@{}lll@{}}
\toprule\noalign{}
Family & Architecture A & Architecture B \\
\midrule\noalign{}
\endhead
\bottomrule\noalign{}
\endlastfoot
Gompertz & 74.3 \(\pm\) 1.4 & 77.4 \(\pm\) 1.3 \\
Logistic & 78.2 \(\pm\) 1.3 & 79.1 \(\pm\) 1.3 \\
Hyperbolic tangent & 76.0 \(\pm\) 1.4 & 77.7 \(\pm\) 1.3 \\
Piecewise-linear breakpoint & 77.8 \(\pm\) 1.3 & 80.4 \(\pm\) 1.3 \\
\end{longtable}
}

\begin{bullets}

\textbf{The inferential consequence of trajectory family is
architecture-dependent: exactly zero under the covariance
channel, modest but non-zero under faithful mean moderation.}

\begin{itemize}\setlength{\itemsep}{2pt}\setlength{\parskip}{0pt}
\item Architecture B: with matched seeds the four families
  return identical rejection decisions, so family is inert by
  construction; the production-seed spread is sampling noise.
\item Architecture A: families span 0.743 to 0.782, a 0.039
  spread; Gompertz (the default) is the lowest-power family,
  roughly two MCSE below logistic.
\item The pre-registered hypothesis that the piecewise-linear
  breakpoint would behave qualitatively differently is not
  supported; it tracks the logistic, not the Gompertz.
\end{itemize}

\end{bullets}

\begin{rgt}
Lorem ipsum dolor sit amet, consectetur adipiscing elit, sed do eiusmod tempor incididunt ut labore et dolore magna aliqua.

\end{rgt}

\begin{orig}
The two architectures behave qualitatively differently. Under
Architecture B (MVN differential correlation), the
biomarker-treatment interaction is carried entirely by the
covariance structure, which is independent of the BR mean curve;
consequently the trajectory family, which sets that mean curve,
cannot alter the interaction test. We confirm this directly: in
a matched-seed check the four families return bit-for-bit
identical rejection decisions, so the production-seed cell spread
(\(0.774\) to \(0.804\)) is sampling noise rather than a family
effect. Under Architecture A (direct mean moderation),
implemented in the faithful multiplicative form in which the
biomarker-dependent component of the treatment effect rides the
response trajectory, the families separate: power ranges from
\(0.743\) (Gompertz) through \(0.760\) (hyperbolic-tangent) and
\(0.778\) (piecewise-linear) to \(0.782\) (logistic), a \(0.039\)
spread. The Gompertz-logistic gap of \(0.039\) is approximately
two standard errors of the difference (SE \(\approx 0.019\)), and
the Gompertz default is the lowest-power, that is the mildly
conservative, family of the four. The pre-registered hypothesis
that the piecewise-linear breakpoint family would behave
qualitatively differently from the smooth families is not
supported in the direction anticipated: the breakpoint family
tracks the logistic and the hyperbolic-tangent, not the Gompertz
default.

\end{orig}

\begin{bullets}

\textbf{The earlier flat-moderation result, which reported no
family effect under either architecture, was an artefact of the
interaction encoding, not a property of trajectory family.}

\begin{itemize}\setlength{\itemsep}{2pt}\setlength{\parskip}{0pt}
\item The prototype applied the biomarker moderation as a flat
  shift, independent of the response trajectory, so family could
  not affect the interaction under either architecture.
\item Once the moderation is trajectory-scaled per the paper-01
  Architecture-A model, the mean-moderation families separate.
\item The covariance-channel invariance is unchanged, because
  that channel is genuinely mean-independent.
\end{itemize}

\end{bullets}

\begin{rgt}
Lorem ipsum dolor sit amet, consectetur adipiscing elit, sed do eiusmod tempor incididunt ut labore et dolore magna aliqua.

\end{rgt}

\begin{orig}
An earlier draft of this manuscript reported essentially no
family effect under either architecture. That result is now
understood to have been an artefact of the interaction encoding
rather than a property of trajectory family. In that draft the
Architecture-A biomarker moderation was applied as a flat
per-timepoint shift of magnitude \(\beta_{bm}\sigma_{BR}\),
independent of the response trajectory, so that the trajectory
family, which enters only through the BR mean curve, could not
affect the bm:Dbc interaction under Architecture A any more than
under the mean-independent Architecture B. Implementing
Architecture A in its faithful multiplicative form (section 4),
in which the biomarker-dependent component scales with the
response trajectory in keeping with the companion manuscript at
\texttt{analysis/report/01-dgp-mean-moderation-vs-mvn/}, the
mean-moderation families separate at the few-percent level while
the covariance-channel invariance under Architecture B is
preserved, the latter being a genuine mean-independence property
rather than an artefact.

\end{orig}

\subsection{Type I error across response-curve families}\label{type-i-error-across-response-curve-families}

Type I error at \(c_{bm} = 0\), percent \(\pm\) MCSE:

{\def\LTcaptype{none} % do not increment counter
\begin{longtable}[]{@{}lll@{}}
\toprule\noalign{}
Family & Architecture A & Architecture B \\
\midrule\noalign{}
\endhead
\bottomrule\noalign{}
\endlastfoot
Gompertz & 1.9 \(\pm\) 0.4 & 2.6 \(\pm\) 0.5 \\
Logistic & 3.8 \(\pm\) 0.6 & 2.2 \(\pm\) 0.5 \\
Hyperbolic tangent & 2.7 \(\pm\) 0.5 & 2.6 \(\pm\) 0.5 \\
Piecewise-linear breakpoint & 2.8 \(\pm\) 0.5 & 2.7 \(\pm\) 0.5 \\
\end{longtable}
}

\begin{bullets}

\textbf{Type I error is slightly conservative across all four
families, consistent with corCAR1 small-sample
degree-of-freedom accounting.}

\begin{itemize}\setlength{\itemsep}{2pt}\setlength{\parskip}{0pt}
\item All eight null cells fall in $[0.019, 0.038]$; no
  family-specific inflation is observed.
\item At the null the moderation term vanishes, so the
  trajectory-scaling of Architecture A is inactive and the
  null cells are governed by the same mechanism as before.
\item The conservatism is shared across families, implicating
  the corCAR1 d.f. mechanism rather than trajectory shape.
\end{itemize}

\end{bullets}

\begin{rgt}
Lorem ipsum dolor sit amet, consectetur adipiscing elit, sed do eiusmod tempor incididunt ut labore et dolore magna aliqua.

\end{rgt}

\begin{orig}
All eight null cells produce empirical Type I error in
\([0.019, 0.038]\), slightly conservative relative to the nominal
0.05 and consistent with the corCAR1 small-sample
degree-of-freedom accounting at the prazosin-calibrated \(N = 35\)
and 8-timepoint short serial series. At the null the biomarker
moderation term is zero, so the trajectory-scaling that
distinguishes the Architecture-A families under the alternative
is inactive here, and the null cells differ across families only
by sampling noise. There is no family-specific inflation; the
highest cell, logistic under Architecture A at 0.038, is within
two MCSE of the others. The conservatism is shared across all
four families, supporting the interpretation that it reflects
the corCAR1 d.f. mechanism rather than a trajectory-shape
artefact, and it has appeared consistently across pmsimstats
simulation runs at the reference parameters.

\end{orig}

\subsection{Bias of the biomarker-treatment interaction estimate}\label{bias-of-the-biomarker-treatment-interaction-estimate}

\begin{bullets}

\textbf{Bias of the interaction estimate is negligible and
family-invariant under both architectures, so the
Architecture-A power separation is a precision effect, not a
bias effect.}

\begin{itemize}\setlength{\itemsep}{2pt}\setlength{\parskip}{0pt}
\item Mean $\hat\beta_{bm:D_{bc}}$ varies from $-0.230$ to
  $-0.239$ across the eight alternative cells, a 0.009
  spread within Monte Carlo noise.
\item Empirical SE is essentially constant across families
  (0.073 to 0.083).
\item The matched-anchor calibration holds the estimand fixed
  across families, so the power spread reflects the temporal
  profile of the signal rather than its magnitude or bias.
\end{itemize}

\end{bullets}

\begin{rgt}
Lorem ipsum dolor sit amet, consectetur adipiscing elit, sed do eiusmod tempor incididunt ut labore et dolore magna aliqua.

\end{rgt}

\begin{orig}
Across the eight alternative cells the mean estimated
\(\hat\beta_{bm:D_{bc}}\) at \(c_{bm} = 0.45\) varies over a narrow
range (roughly \(-0.230\) to \(-0.239\)), a cell-to-cell spread of
about 0.009 that is itself within the Monte Carlo noise of the
estimator and shows no family-specific drift; the Gompertz and
logistic cells that differ most in power have nearly identical
mean estimates. The empirical SE of \(\hat\beta_{bm:D_{bc}}\) is
essentially constant across families (SD in the range 0.073 to
0.083 across cells) and roughly equal between architectures.
That the Architecture-A power separation appears with negligible
bias and near-constant empirical SE indicates it is a precision
effect operating through the test statistic's sensitivity to the
temporal profile of the signal, not a shift in the estimand or
its sampling dispersion. The matched-anchor calibration holds
the marginal effect size fixed across families by construction,
so the residual family differences are differences in the
temporal shape of the moderation, which is precisely the axis
the design isolates.

\end{orig}

\section{Discussion}\label{discussion}

\subsection{An architecture-dependent finding}\label{an-architecture-dependent-finding}

\begin{bullets}

\textbf{The inferential consequence of trajectory family depends
on how the biomarker-treatment interaction is encoded: zero
under the covariance channel, modest under mean moderation.}

\begin{itemize}\setlength{\itemsep}{2pt}\setlength{\parskip}{0pt}
\item Under Architecture B (covariance channel) the family is
  inert by construction; the interaction test is mean-independent.
\item Under Architecture A (mean moderation) the four families
  separate over a 0.039 power range, with the Gompertz default
  the lowest-power, that is mildly conservative, family.
\item The effect is modest (about two MCSE at $N = 35$,
  $c_{bm} = 0.45$) and should not be overstated, but it is real
  and direction-consistent.
\end{itemize}

\end{bullets}

\begin{rgt}
Lorem ipsum dolor sit amet, consectetur adipiscing elit, sed do eiusmod tempor incididunt ut labore et dolore magna aliqua.

\end{rgt}

\begin{orig}
The central finding is that the inferential consequence of
trajectory-family choice is architecture-dependent. At the
prazosin-calibrated reference cell (Hybrid OL+BDC, \(N = 35\),
\(c_{bm} = 0.45\), \(t_{1/2} = 0\), matched-anchor calibration),
the four families span the qualitative range of plausible
templates for monotone-saturating treatment-response
trajectories: asymmetric (Gompertz), symmetric smooth (logistic
and hyperbolic-tangent), and non-smooth (piecewise-linear
breakpoint). Under Architecture B, where the interaction is
carried by the BM-BR covariance and is therefore independent of
the BR mean curve, the family is inert by construction. Under
Architecture A, where the interaction is a mean effect that, in
its faithful multiplicative form, scales with the response
trajectory, the families separate over a 0.039 power range and
the Gompertz default is the lowest-power, that is the mildly
conservative, family. We would not wish to overstate the
Architecture-A effect; it is modest, on the order of two Monte
Carlo standard errors at this sample size and effect size, and
the conditions under which it would grow or shrink are discussed
in section 6.3. But it is real and direction-consistent, and it
overturns the apparent null reported in the earlier flat-
moderation draft.

\end{orig}

\begin{bullets}

\textbf{The pre-registered piecewise-linear hypothesis is not
supported in the anticipated direction; the breakpoint family
tracks the smooth alternatives, not the Gompertz default.}

\begin{itemize}\setlength{\itemsep}{2pt}\setlength{\parskip}{0pt}
\item The non-differentiable breakpoint carries no distinctive
  power signature at $N = 35$ and $c_{bm} = 0.45$.
\item Under Architecture A the breakpoint family sits with the
  logistic and hyperbolic-tangent, above the Gompertz default.
\item Smoothness per se is not the operative axis; the temporal
  profile of the moderation is.
\end{itemize}

\end{bullets}

\begin{rgt}
Lorem ipsum dolor sit amet, consectetur adipiscing elit, sed do eiusmod tempor incididunt ut labore et dolore magna aliqua.

\end{rgt}

\begin{orig}
The pre-registered hypothesis that the piecewise-linear
breakpoint family, with its non-differentiable feature, would
behave qualitatively differently from the smooth families is not
supported in the direction anticipated. The breakpoint family
carries no distinctive power signature attributable to its
non-smoothness; under Architecture A it sits with the logistic
and hyperbolic-tangent families, above the Gompertz default, and
under Architecture B it is inert like the others. The operative
axis is therefore not smoothness but the temporal profile of the
biomarker-moderated signal: families whose response rises more
gradually across the on-drug timepoints, and so carry more of
the moderated signal at the later, more informative timepoints,
yield marginally higher power than the front-loaded Gompertz
curve, which reaches near-asymptote early.

\end{orig}

\subsection{Why the architectures differ}\label{why-the-architectures-differ}

\begin{bullets}

\textbf{Two mechanisms explain the contrast: exact mean-
independence of the covariance channel, and the residual
profile-sensitivity of the trajectory-scaled mean channel.}

\begin{itemize}\setlength{\itemsep}{2pt}\setlength{\parskip}{0pt}
\item Under Architecture B the biomarker-treatment information
  lives in the covariance and is mathematically independent of
  the trajectory mean, so the family cannot move the test.
\item Under Architecture A the moderation rides the trajectory,
  so the temporal profile of the signal, not just its anchored
  magnitude, enters the Wald statistic.
\item Matched-anchor calibration leaves only profile differences,
  which the time-averaging Wald test damps but does not erase.
\end{itemize}

\end{bullets}

\begin{rgt}
Lorem ipsum dolor sit amet, consectetur adipiscing elit, sed do eiusmod tempor incididunt ut labore et dolore magna aliqua.

\end{rgt}

\begin{orig}
We see two mechanisms that, taken together, account for the
architecture contrast, and we shall work through each in turn.
The first is the exact mean-independence of the covariance
channel under Architecture B; the second is the residual
profile-sensitivity of the trajectory-scaled mean channel under
Architecture A, which the matched-anchor calibration and the
time-averaging of the Wald test damp but do not eliminate.

\end{orig}

\begin{bullets}

\textbf{Matched-anchor calibration ensures that residual family
differences are in shape only, not in amplitude, and those
shape differences are small relative to between-replicate
variation.}

\begin{itemize}\setlength{\itemsep}{2pt}\setlength{\parskip}{0pt}
\item Families matched at asymptote and $t = 5$ anchor.
\item Within-trajectory shape variation is small relative to
  the between-replicate variance of the bm:Dbc estimator.
\item The calibration design isolates shape effects from
  amplitude effects by construction.
\end{itemize}

\end{bullets}

\begin{rgt}
Lorem ipsum dolor sit amet, consectetur adipiscing elit, sed do eiusmod tempor incididunt ut labore et dolore magna aliqua.

\end{rgt}

\begin{orig}
The first mechanism damps the Architecture-A effect: calibration
removes the amplitude differences, leaving only profile
differences. The four families are matched at the asymptotic
maximum and at a single anchor point (\(t = 5\) weeks), and the
trajectory-scaled moderation is normalised to a common mean
on-drug magnitude. At matched marginal effect size, the
residual within-trajectory variation in shape is small relative
to the between-replicate variation in the bm:Dbc estimator. The
calibration constraint is intentional in the simulation design:
it isolates trajectory-shape effects from amplitude effects, so
that the family-specific power differences we observe under
Architecture A can be attributed to the temporal profile of the
moderation alone, and explains why that effect is on the order
of a few percentage points rather than larger.

\end{orig}

\begin{bullets}

\textbf{The bm:Dbc Wald test averages over on-drug timepoints,
smoothing out local trajectory-shape differences that would be
visible in point-in-time comparisons.}

\begin{itemize}\setlength{\itemsep}{2pt}\setlength{\parskip}{0pt}
\item The continuous $D_{bc}$ predictor integrates trajectory
  contributions across all on-drug timepoints.
\item Local shape differences are averaged within the
  integration and do not accumulate to detectable power
  differences.
\item This temporal averaging damps, but does not erase, the
  profile differences that drive the Architecture-A spread.
\end{itemize}

\end{bullets}

\begin{rgt}
Lorem ipsum dolor sit amet, consectetur adipiscing elit, sed do eiusmod tempor incididunt ut labore et dolore magna aliqua.

\end{rgt}

\begin{orig}
The second mechanism, also damping, is that the bm:Dbc Wald test
averages over time. The continuous \(D_{bc}\) predictor multiplied
by the biomarker integrates the trajectory's contribution across
all on-drug timepoints, with weights determined by the
predictor's exponential-decay structure. Trajectory-shape
differences (where one family rises faster than another between
timepoints) are partly averaged within the integration, so the
Wald test is relatively, though not perfectly, insensitive to
profile features. Under Architecture A the averaging is
incomplete: the residual profile difference survives into the
test statistic and produces the observed few-percent power
spread. Under Architecture B there is no profile term to survive,
because the interaction does not enter the mean at all.

\end{orig}

\begin{bullets}

\textbf{Under Architecture B the interaction is in the
covariance structure and is exactly independent of trajectory
shape, which we verify by a matched-seed check.}

\begin{itemize}\setlength{\itemsep}{2pt}\setlength{\parskip}{0pt}
\item Architecture B encodes information in
  $\partial \mathrm{Cor}(BM, BR)/\partial D_{bc}$, not in
  the trajectory mean.
\item Holding the seed fixed, the four families return identical
  rejection decisions, so the invariance is exact, not merely
  within Monte Carlo noise.
\item This is the mechanism behind the exact Architecture-B
  inertness, in contrast to the merely damped Architecture-A
  effect.
\end{itemize}

\end{bullets}

\begin{rgt}
Lorem ipsum dolor sit amet, consectetur adipiscing elit, sed do eiusmod tempor incididunt ut labore et dolore magna aliqua.

\end{rgt}

\begin{orig}
Mean-independence applies specifically to Architecture B: the
covariance moderation is upstream of trajectory shape. Under MVN
differential correlation, the biomarker-treatment information
lives in the \(\partial \mathrm{Cor}(BM, BR)/\partial D_{bc}\)
structure, which is independent of the trajectory family that BR
follows; the family enters only the BR mean, to which the
interaction test is invariant. This is not a small-effect
statement but an exact one, and we verify it directly: holding
the random-number seed fixed across families, the four
Architecture-B cells return bit-for-bit identical rejection
decisions. The production-seed spread under Architecture B is
therefore sampling noise, whereas the Architecture-A spread is a
genuine, if modest, profile effect.

\end{orig}

\begin{bullets}

\textbf{The two mechanisms together place the modified Gompertz
default as inferentially free under the covariance channel and
mildly conservative under the mean channel.}

\begin{itemize}\setlength{\itemsep}{2pt}\setlength{\parskip}{0pt}
\item Under Architecture B the default is robust by exact
  construction; trajectory family cannot matter.
\item Under Architecture A the default is the lowest-power
  family, so it is a safe, slightly conservative, choice rather
  than an optimal one.
\item The following section identifies conditions that would
  amplify the Architecture-A effect.
\end{itemize}

\end{bullets}

\begin{rgt}
Lorem ipsum dolor sit amet, consectetur adipiscing elit, sed do eiusmod tempor incididunt ut labore et dolore magna aliqua.

\end{rgt}

\begin{orig}
The combined effect is that trajectory-family choice is
inferentially free under the covariance channel and carries a
modest, mildly conservative, cost under the mean channel. The
framework's modified Gompertz default is robust by exact construction
under Architecture B, and under Architecture A it is the
lowest-power of the four families, which makes it a safe rather
than an optimal default: a trial designer who adopts it under a
mean-moderation DGP will, if anything, slightly under-state
achievable power. We turn next to the conditions under which the
Architecture-A effect would be expected to grow.

\end{orig}

\subsection{Conditions that would change the answer}\label{conditions-that-would-change-the-answer}

\begin{bullets}

\textbf{Four conditions outside the present 16-cell factorial
would plausibly amplify the Architecture-A family effect; none
would activate an effect under the mean-independent Architecture
B.}

\begin{itemize}\setlength{\itemsep}{2pt}\setlength{\parskip}{0pt}
\item Smaller $N$, different $c_{bm}$, non-zero carryover
  ($t_{1/2} > 0$), and non-monotone or biphasic shape
  envelopes.
\item None of these conditions is exercised in the present
  focused factorial.
\item Each is a natural axis for the 144-cell follow-up grid;
  the biphasic case is the only one expected to reach
  Architecture B.
\end{itemize}

\end{bullets}

\begin{rgt}
Lorem ipsum dolor sit amet, consectetur adipiscing elit, sed do eiusmod tempor incididunt ut labore et dolore magna aliqua.

\end{rgt}

\begin{orig}
Four conditions, none exercised in the present 16-cell
factorial, would plausibly amplify the modest Architecture-A
family effect, and each merits a brief comment. Because
Architecture B is mean-independent, the first three conditions
(sample size, effect size, carryover) bear on Architecture A
only; the fourth, non-monotone shape, is the only one that could
in principle reach Architecture B, and then only by changing the
covariance construction rather than the mean.

\end{orig}

\begin{bullets}

\textbf{Smaller sample size might amplify local-shape
sensitivity by concentrating temporal information at fewer
timepoints.}

\begin{itemize}\setlength{\itemsep}{2pt}\setlength{\parskip}{0pt}
\item At $N = 35$, the test integrates over a moderately
  large within-subject sample.
\item With fewer timepoints, local shape differences between
  families would be more visible rather than averaged.
\item This condition is the most proximate candidate for
  amplifying the Architecture-A effect beyond a few percent.
\end{itemize}

\end{bullets}

\begin{rgt}
Lorem ipsum dolor sit amet, consectetur adipiscing elit, sed do eiusmod tempor incididunt ut labore et dolore magna aliqua.

\end{rgt}

\begin{orig}
First, smaller sample size. At \(N = 35\) the bm:Dbc test
integrates over a moderately large within-subject sample;
smaller \(N\) might amplify the local-shape sensitivity by
concentrating the temporal information at fewer timepoints
where the families differ more visibly.

\end{orig}

\begin{bullets}

\textbf{At smaller $c_{bm}$ near the 50\% power boundary,
family differences might surface via different finite-sample
sampling distributions.}

\begin{itemize}\setlength{\itemsep}{2pt}\setlength{\parskip}{0pt}
\item The $c_{bm} = 0.45$ cell sits in a moderate-power
  regime with comparable signal-to-noise across families.
\item Pushing to $c_{bm} = 0.30$ (power near 0.5) may reveal
  family-specific finite-sample sampling differences.
\item The tight clustering at 0.45 may be a local feature of
  the effect-size space.
\end{itemize}

\end{bullets}

\begin{rgt}
Lorem ipsum dolor sit amet, consectetur adipiscing elit, sed do eiusmod tempor incididunt ut labore et dolore magna aliqua.

\end{rgt}

\begin{orig}
Second, larger or smaller \(c_{bm}\). The \(c_{bm} = 0.45\)
reference cell sits in a power regime where families have
comparable signal-to-noise. Pushing \(c_{bm}\) to 0.30 (where
power lies near 0.5) might surface family differences via
different finite-sample sampling distributions, and it is
possible that the tight clustering we observe is a feature of
this particular region of the effect-size space.

\end{orig}

\begin{bullets}

\textbf{Non-zero carryover introduces an interaction between
trajectory shape and the off-drug exponential decay that is
suppressed in the present $t_{1/2} = 0$ slice.}

\begin{itemize}\setlength{\itemsep}{2pt}\setlength{\parskip}{0pt}
\item With $t_{1/2} > 0$, integration weights interact with
  the off-drug decay, potentially differentiating families.
\item The present simulation sets $t_{1/2} = 0$ to isolate
  trajectory shape from carryover effects.
\item This interaction is the mechanistic pathway by which
  carryover could make trajectory family consequential.
\end{itemize}

\end{bullets}

\begin{rgt}
Lorem ipsum dolor sit amet, consectetur adipiscing elit, sed do eiusmod tempor incididunt ut labore et dolore magna aliqua.

\end{rgt}

\begin{orig}
Third, carryover-induced correlation decay. With \(t_{1/2} > 0\),
the bm:Dbc-mediated information channel is eroded differently
for different trajectory shapes, because the integration weights
interact with the off-drug exponential decay. The \(t_{1/2} = 0\)
slice exercised here suppresses this interaction entirely.

\end{orig}

\begin{bullets}

\textbf{Non-monotone or biphasic trajectory families would
produce genuinely different bm:Dbc information content to
which the Wald test would almost certainly respond.}

\begin{itemize}\setlength{\itemsep}{2pt}\setlength{\parskip}{0pt}
\item All four present families share a monotone-saturating
  envelope; the small-effect finding is conditional on this
  property.
\item Biphasic or oscillating families have qualitatively
  different information structure.
\item Extension to non-monotone families is outside the scope
  of the present study.
\end{itemize}

\end{bullets}

\begin{rgt}
Lorem ipsum dolor sit amet, consectetur adipiscing elit, sed do eiusmod tempor incididunt ut labore et dolore magna aliqua.

\end{rgt}

\begin{orig}
Fourth, non-monotone or biphasic shape envelopes. The four
families we have examined share a monotone-saturating envelope.
Families that peak and decay (biphasic) or oscillate would
produce genuinely different trajectories with different bm:Dbc
information content, and the Wald test would almost certainly
respond to those differences.

\end{orig}

\begin{bullets}

\textbf{The 144-cell extended grid is the canonical follow-up
that would close the trajectory-family question at publication-
adequate rigour.}

\begin{itemize}\setlength{\itemsep}{2pt}\setlength{\parskip}{0pt}
\item The grid crosses $N$, $c_{bm}$, $t_{1/2}$, and family.
\item Each of the four conditions above maps to a dimension of
  this grid.
\item The present focused factorial is a necessary precursor;
  the extended grid is the production-grade completion.
\end{itemize}

\end{bullets}

\begin{rgt}
Lorem ipsum dolor sit amet, consectetur adipiscing elit, sed do eiusmod tempor incididunt ut labore et dolore magna aliqua.

\end{rgt}

\begin{orig}
Each condition is the natural axis of a follow-up sensitivity
analysis. The 144-cell grid anticipated in the abstract
(\$N \times c\_\{bm\} \times t\_\{1/2\} \times \$ family) is the
canonical follow-up, and we recommend it as the production-grade
extension that would close out the trajectory-family question
at a level of rigour adequate for the methodological literature.

\end{orig}

\subsection{The prazosin/PTSD reference cell}\label{the-prazosinptsd-reference-cell}

\begin{bullets}

\textbf{For the prazosin-PTSD application, trajectory family is a
low-consequence sensitivity axis relative to carryover and
DGP-architecture, but not a free one under mean moderation.}

\begin{itemize}\setlength{\itemsep}{2pt}\setlength{\parskip}{0pt}
\item Sensitivity to the modified Gompertz function is small relative
  to carryover specification and DGP-architecture sensitivity.
\item Under a covariance-channel DGP the modified Gompertz choice is
  inferentially free; under mean moderation it is mildly
  conservative, costing a few percentage points of power.
\item Sensitivity-analysis resources are better spent on
  carryover and DGP-architecture axes with larger inferential
  consequences, while still reporting one family comparison
  under the mean-moderation DGP.
\end{itemize}

\end{bullets}

\begin{rgt}
Lorem ipsum dolor sit amet, consectetur adipiscing elit, sed do eiusmod tempor incididunt ut labore et dolore magna aliqua.

\end{rgt}

\begin{orig}
For the motivating prazosin-PTSD application that drives the
pmsimstats research program, the present result implies that
sensitivity to the modified Gompertz assumption is small relative to
sensitivity to the carryover specification (the subject of the
companion manuscript at
\texttt{analysis/report/02-carryover-sensitivity/}) and to
the DGP architecture (the subject of
\texttt{analysis/report/01-dgp-mean-moderation-vs-mvn/}). Of the
three principal axes of sensitivity uncertainty in the
pmsimstats simulation framework, trajectory family choice is, at
the reference parameters, the least consequential by a
substantial margin. It is not, however, entirely free: under the
mean-moderation DGP the modified Gompertz default is the lowest-power
family by a few percentage points, so a trial designer who
adopts it is making a mildly conservative choice. Under the
covariance-channel DGP the choice is inferentially free. We
would suggest that sensitivity-analysis attention be directed
chiefly to the carryover and DGP-architecture axes, where the
inferential consequences are larger, while reporting at least
one Gompertz-vs-alternative comparison whenever the operative
DGP is mean moderation.

\end{orig}

\section{Conclusions}\label{conclusions}

\begin{bullets}

\textbf{The modified Gompertz is a defensible template for
monotone-saturating N-of-1 trajectories; the inferential cost of
the choice is architecture-dependent.}

\begin{itemize}\setlength{\itemsep}{2pt}\setlength{\parskip}{0pt}
\item Under the covariance channel (Architecture B) the four
  families are exactly equivalent; under mean moderation
  (Architecture A) they span a 0.039 power range at MCSE
  0.013--0.014.
\item The Gompertz default is the lowest-power family under mean
  moderation, that is mildly conservative, not optimal.
\item The piecewise-linear breakpoint hypothesis is not
  supported; the breakpoint family tracks the smooth alternatives.
\end{itemize}

\end{bullets}

\begin{rgt}
Lorem ipsum dolor sit amet, consectetur adipiscing elit, sed do eiusmod tempor incididunt ut labore et dolore magna aliqua.

\end{rgt}

\begin{orig}
First, the modified Gompertz function is a defensible
parametric template for monotone-saturating treatment-response
trajectories in N-of-1 simulation, with an architecture-dependent
qualification. At the prazosin/PTSD reference cell with
matched-anchor calibration, the four-family comparison (Gompertz,
symmetric logistic, hyperbolic-tangent, piecewise-linear
breakpoint) produces exactly equivalent inference under the
covariance-channel architecture, where the interaction is
mean-independent, and a modest 0.039 power spread under the
mean-moderation architecture, where the biomarker moderation
rides the response trajectory. In the latter case the Gompertz
default is the lowest-power of the four families, which makes it
a mildly conservative rather than an optimal default. The
pre-registered hypothesis that the piecewise-linear breakpoint
family would behave qualitatively differently from the smooth
families is not supported in the direction anticipated.

\end{orig}

\begin{bullets}

\textbf{The finding implies the Wald test is robust to
trajectory family under the covariance channel and mildly
sensitive to it under mean moderation; boundary conditions are
deferred to the 144-cell follow-up.}

\begin{itemize}\setlength{\itemsep}{2pt}\setlength{\parskip}{0pt}
\item Robustness is exact under Architecture B and approximate
  (few-percent) under Architecture A in the evaluated regime.
\item At what sample sizes, effect sizes, or carryover levels
  the Architecture-A effect grows is an open question.
\item Sensitivity-analysis priority should go to carryover and
  DGP-architecture axes with larger demonstrated consequences.
\end{itemize}

\end{bullets}

\begin{rgt}
Lorem ipsum dolor sit amet, consectetur adipiscing elit, sed do eiusmod tempor incididunt ut labore et dolore magna aliqua.

\end{rgt}

\begin{orig}
Second, the finding is methodologically informative precisely
because it is architecture-dependent. It implies that the bm:Dbc
Wald test is exactly robust to trajectory-family choice when the
interaction is encoded in the covariance, and only mildly
sensitive to it, at the level of a few percentage points of
power, when the interaction is a trajectory-scaled mean effect.
We cannot yet establish at what sample sizes, effect sizes, or
carryover levels the Architecture-A effect would grow; that
question is deferred to the 144-cell follow-up. Trial designers
using the pmsimstats framework can adopt the modified Gompertz default
with confidence under a covariance-channel DGP, and with the
understanding that it is mildly conservative under a
mean-moderation DGP, while directing sensitivity-analysis
attention to the carryover and DGP-architecture axes that have
larger inferential consequences.

\end{orig}

\begin{bullets}

\textbf{The 144-cell extended grid is the work needed to close
the trajectory-family question at publication-adequate rigour.}

\begin{itemize}\setlength{\itemsep}{2pt}\setlength{\parskip}{0pt}
\item The present 16-cell factorial addresses the central
  question with statistical clarity.
\item The extended grid would map boundary conditions at which
  trajectory family begins to matter.
\item This follow-up is recommended as the production-grade
  completion of the trajectory-family evaluation.
\end{itemize}

\end{bullets}

\begin{rgt}
Lorem ipsum dolor sit amet, consectetur adipiscing elit, sed do eiusmod tempor incididunt ut labore et dolore magna aliqua.

\end{rgt}

\begin{orig}
Third, the 144-cell extended grid (sample sizes, effect sizes,
DGP-side carryover) anticipated in the original abstract is the
natural follow-up, and we regard it as the work that would be
needed to close the trajectory-family question at a level of
rigour adequate for the methodological literature. The present
focused 16-cell factorial addresses the central scientific
question with statistical clarity; the extended grid would map
the boundary conditions at which trajectory-family choice
begins to matter.

\end{orig}

\section{Conflict of interest}\label{conflict-of-interest}

The authors declare no conflicts of interest.

\section{Funding}\label{funding}

This work received no specific funding.

\section{Data availability statement}\label{data-availability-statement}

\begin{bullets}

\textbf{Data and code are openly available in the
pmsimstats-ng repository under versioned standard paths.}

\begin{itemize}\setlength{\itemsep}{2pt}\setlength{\parskip}{0pt}
\item Per-replicate Monte Carlo outputs and ADEMP cell-level
  summaries are at \texttt{analysis/data/}.
\item Simulation drivers are at \texttt{analysis/scripts/}.
\item Exact manuscript paths are listed in the Reproducibility
  section.
\end{itemize}

\end{bullets}

\begin{rgt}
Lorem ipsum dolor sit amet, consectetur adipiscing elit, sed do eiusmod tempor incididunt ut labore et dolore magna aliqua.

\end{rgt}

\begin{orig}
The data and code that support the findings of this study are
openly available in the pmsimstats-ng project repository.
Per-replicate Monte Carlo outputs and ADEMP-compliant
cell-level summaries are versioned under
\texttt{analysis/data/}; simulation drivers are at
\texttt{analysis/scripts/}. Exact paths for this manuscript
are listed in the Reproducibility section below.

\end{orig}

\section{Reproducibility}\label{reproducibility}

\begin{bullets}

\textbf{The manuscript is fully reproducible inside the
zzcollab Docker container with the renv.lock package set.}

\begin{itemize}\setlength{\itemsep}{2pt}\setlength{\parskip}{0pt}
\item Container image hash in \texttt{Dockerfile}; package
  set in \texttt{renv.lock}.
\item Principal packages: \texttt{nlme}, \texttt{parallel},
  \texttt{data.table}, \texttt{furrr}/\texttt{purrr},
  \texttt{rmarkdown}/\texttt{bookdown}.
\item Rendered with R 4.4.x.
\end{itemize}

\end{bullets}

\begin{rgt}
Lorem ipsum dolor sit amet, consectetur adipiscing elit, sed do eiusmod tempor incididunt ut labore et dolore magna aliqua.

\end{rgt}

\begin{orig}
This manuscript was rendered with R 4.4.x inside the project's
zzcollab Docker container (image hash recorded in
\texttt{Dockerfile}; package set captured in
\texttt{renv.lock}). Principal packages used by the simulation
pipeline are \texttt{nlme} (continuous-time linear-mixed
analysis with \texttt{corCAR1} residual correlation),
\texttt{parallel} (8-core \texttt{mclapply} for
parameter-grid sweeps), \texttt{data.table} (the
\texttt{original-extended} simulation collection),
\texttt{furrr} and \texttt{purrr} (the \texttt{tidyverse}
simulation collection), and \texttt{rmarkdown} plus
\texttt{bookdown} for the manuscript build.

\end{orig}

\begin{bullets}

\textbf{Seeds, checkpoints, and driver paths are fully
specified for exact replication of the simulation results.}

\begin{itemize}\setlength{\itemsep}{2pt}\setlength{\parskip}{0pt}
\item Master seed \texttt{set.seed(20260617)}; per-cell seed
  offsets give independent but reproducible replicate streams.
\item Outputs checkpointed at
  \texttt{analysis/data/quick-sim/07-gompertz-faithful-replicates.rds};
  ADEMP summaries at \texttt{07-gompertz-faithful-summary.txt}.
\item Driver \texttt{02-faithful-trajectory-sweep.R}; pre-registered
  ADEMP plan at
  \texttt{analysis/scripts/gompertz-evaluation/00-ademp-pre-registration.md}.
\end{itemize}

\end{bullets}

\begin{rgt}
Lorem ipsum dolor sit amet, consectetur adipiscing elit, sed do eiusmod tempor incididunt ut labore et dolore magna aliqua.

\end{rgt}

\begin{orig}
The simulation driver
\texttt{analysis/scripts/gompertz-evaluation/02-faithful-trajectory-sweep.R}
sets \texttt{set.seed(20260617)} and assigns each of the 16 cells
a distinct seed offset, so replicate streams are independent
across cells but reproducible. The trajectory family is routed
through the data-generating process via the \texttt{br\_family}
argument to \texttt{buildSigma}, and the Architecture-A
moderation is trajectory-scaled via the
\texttt{moderation\_scaling = "trajectory"} argument to
\texttt{generateData} (both default to the original behavior, so
all other pmsimstats runs are unchanged). Per-replicate
simulation outputs are checkpointed at
\texttt{analysis/data/quick-sim/07-gompertz-faithful-replicates.rds};
Morris ADEMP cell-level summaries are at the companion
\texttt{07-gompertz-faithful-summary.txt}. The pre-registered
ADEMP plan is at
\texttt{analysis/scripts/gompertz-evaluation/00-ademp-pre-registration.md},
and deviations from it are recorded in
\texttt{02-deviations-log.md}. The manuscript source,
paper-specific bibliography (\texttt{references.bib}), and SIM
CSL file are at
\texttt{analysis/report/07-gompertz-evaluation/}.

\end{orig}

\section{References}\label{references}

\protect\phantomsection\label{refs}
\begin{CSLReferences}{0}{1}
\bibitem[\citeproctext]{ref-hendrickson2020}
\CSLLeftMargin{1. }%
\CSLRightInline{Hendrickson RC, Thomas RG, Schork NJ, Raskind MA. Optimizing aggregated {N}-of-1 trial designs for predictive biomarker validation: Statistical methods and practical considerations. \emph{Frontiers in Digital Health}. 2020;2:13. doi:\href{https://doi.org/10.3389/fdgth.2020.00013}{10.3389/fdgth.2020.00013}}

\bibitem[\citeproctext]{ref-laird1964dynamics}
\CSLLeftMargin{2. }%
\CSLRightInline{Laird AK. Dynamics of tumor growth. \emph{British Journal of Cancer}. 1964;13(3):490-502. doi:\href{https://doi.org/10.1038/bjc.1964.55}{10.1038/bjc.1964.55}}

\bibitem[\citeproctext]{ref-norton1988gompertzian}
\CSLLeftMargin{3. }%
\CSLRightInline{Norton L. A {G}ompertzian model of human breast cancer growth. \emph{Cancer Research}. 1988;48(24 Pt 1):7067-7071.}

\bibitem[\citeproctext]{ref-zwietering1990modeling}
\CSLLeftMargin{4. }%
\CSLRightInline{Zwietering MH, Jongenburger I, Rombouts FM, Riet K van 't. Modeling of the bacterial growth curve. \emph{Applied and Environmental Microbiology}. 1990;56(6):1875-1881. doi:\href{https://doi.org/10.1128/aem.56.6.1875-1881.1990}{10.1128/aem.56.6.1875-1881.1990}}

\bibitem[\citeproctext]{ref-richards1959flexible}
\CSLLeftMargin{5. }%
\CSLRightInline{Richards FJ. A flexible growth function for empirical use. \emph{Journal of Experimental Botany}. 1959;10(2):290-301.}

\bibitem[\citeproctext]{ref-tjorve2010richards}
\CSLLeftMargin{6. }%
\CSLRightInline{Tjørve E, Tjørve KMC. A unified approach to the {R}ichards-model family for use in growth analyses: Why we need only two model forms. \emph{Journal of Theoretical Biology}. 2010;267(3):417-425. doi:\href{https://doi.org/10.1016/j.jtbi.2010.09.008}{10.1016/j.jtbi.2010.09.008}}

\bibitem[\citeproctext]{ref-tjorve2017}
\CSLLeftMargin{7. }%
\CSLRightInline{Tjørve KMC, Tjørve E. The use of {G}ompertz models in growth analyses, and new {G}ompertz-model approach: An addition to the {U}nified-{R}ichards family. \emph{PLOS ONE}. 2017;12(6):e0178691. doi:\href{https://doi.org/10.1371/journal.pone.0178691}{10.1371/journal.pone.0178691}}

\bibitem[\citeproctext]{ref-lillie2011}
\CSLLeftMargin{8. }%
\CSLRightInline{Lillie EO, Patay B, Diamant J, Issell B, Topol EJ, Schork NJ. The n-of-1 clinical trial: The ultimate strategy for individualizing medicine? \emph{Personalized Medicine}. 2011;8(2):161-173. doi:\href{https://doi.org/10.2217/pme.11.7}{10.2217/pme.11.7}}

\bibitem[\citeproctext]{ref-schork2023precision}
\CSLLeftMargin{9. }%
\CSLRightInline{Schork NJ, Beaulieu-Jones B, Liang WS, Smalley S, Goetz LH. Exploring human biology with {N}-of-1 clinical trials. \emph{Cambridge Prisms: Precision Medicine}. 2023;1:e12. doi:\href{https://doi.org/10.1017/pcm.2022.15}{10.1017/pcm.2022.15}}

\bibitem[\citeproctext]{ref-vohra2015cent}
\CSLLeftMargin{10. }%
\CSLRightInline{{Vohra S, Shamseer L, Sampson M, et al.} {CONSORT} extension for reporting {N}-of-1 trials ({CENT}) 2015 statement. \emph{Journal of Clinical Epidemiology}. 2015;76:9-17. doi:\href{https://doi.org/10.1016/j.jclinepi.2015.05.004}{10.1016/j.jclinepi.2015.05.004}}

\bibitem[\citeproctext]{ref-benzekry2014classical}
\CSLLeftMargin{11. }%
\CSLRightInline{Benzekry S, Lamont C, Beheshti A, et al. Classical mathematical models for description and prediction of experimental tumor growth. \emph{PLOS Computational Biology}. 2014;10(8):e1003800. doi:\href{https://doi.org/10.1371/journal.pcbi.1003800}{10.1371/journal.pcbi.1003800}}

\bibitem[\citeproctext]{ref-goutelle2008hill}
\CSLLeftMargin{12. }%
\CSLRightInline{Goutelle S, Maurin M, Rougier F, et al. The {H}ill equation: A review of its capabilities in pharmacological modelling. \emph{Fundamental and Clinical Pharmacology}. 2008;22(6):633-648. doi:\href{https://doi.org/10.1111/j.1472-8206.2008.00633.x}{10.1111/j.1472-8206.2008.00633.x}}

\bibitem[\citeproctext]{ref-dayneka1993indirect}
\CSLLeftMargin{13. }%
\CSLRightInline{Dayneka NL, Garg V, Jusko WJ. Comparison of four basic models of indirect pharmacodynamic responses. \emph{Journal of Pharmacokinetics and Biopharmaceutics}. 1993;21(4):457-478. doi:\href{https://doi.org/10.1007/BF01061691}{10.1007/BF01061691}}

\bibitem[\citeproctext]{ref-vonbertalanffy1957}
\CSLLeftMargin{14. }%
\CSLRightInline{Bertalanffy L von. Quantitative laws in metabolism and growth. \emph{Quarterly Review of Biology}. 1957;32(3):217-231.}

\bibitem[\citeproctext]{ref-morris2019}
\CSLLeftMargin{15. }%
\CSLRightInline{Morris TP, White IR, Crowther MJ. Using simulation studies to evaluate statistical methods. \emph{Statistics in Medicine}. 2019;38(11):2074-2102.}

\bibitem[\citeproctext]{ref-jagadeesan2023sloppiness}
\CSLLeftMargin{16. }%
\CSLRightInline{Jagadeesan P, Raman K, Tangirala AK. Sloppiness: Fundamental study, new formalism and its application in model assessment. \emph{PLOS ONE}. 2023;18(3):e0282609. doi:\href{https://doi.org/10.1371/journal.pone.0282609}{10.1371/journal.pone.0282609}}

\bibitem[\citeproctext]{ref-hill1910}
\CSLLeftMargin{17. }%
\CSLRightInline{Hill AV. The possible effects of the aggregation of the molecules of haemoglobin on its dissociation curves. \emph{Journal of Physiology}. 1910;40:iv-vii.}

\bibitem[\citeproctext]{ref-raskind2013prazosin}
\CSLLeftMargin{18. }%
\CSLRightInline{{Raskind MA, Peterson K, Williams T, et al.} A trial of prazosin for combat trauma {PTSD} with nightmares in active-duty soldiers returned from {I}raq and {A}fghanistan. \emph{American Journal of Psychiatry}. 2013;170(9):1003-1010. doi:\href{https://doi.org/10.1176/appi.ajp.2013.12081133}{10.1176/appi.ajp.2013.12081133}}

\bibitem[\citeproctext]{ref-ebling1996feasibility}
\CSLLeftMargin{19. }%
\CSLRightInline{Ebling WF, Matsumoto Y, Levy G. Feasibility of effect-controlled clinical trials of drugs with pharmacodynamic hysteresis using sparse data. \emph{Pharmaceutical Research}. 1996;13(12):1804-1810. doi:\href{https://doi.org/10.1023/a:1016072806164}{10.1023/a:1016072806164}}

\bibitem[\citeproctext]{ref-triggs1996levodopa}
\CSLLeftMargin{20. }%
\CSLRightInline{{Triggs EJ, Charles BG, Contin M, et al.} Population pharmacokinetics and pharmacodynamics of oral levodopa in parkinsonian patients. \emph{European Journal of Clinical Pharmacology}. 1996;51(1):59-67. doi:\href{https://doi.org/10.1007/s002280050161}{10.1007/s002280050161}}

\bibitem[\citeproctext]{ref-hockenberry2017symptom}
\CSLLeftMargin{21. }%
\CSLRightInline{{Hockenberry MJ, Hooke MC, Rodgers C, et al.} Symptom trajectories in children receiving treatment for leukemia: A latent class growth analysis with multitrajectory modeling. \emph{Journal of Pain and Symptom Management}. 2017;54(1):1-8. doi:\href{https://doi.org/10.1016/j.jpainsymman.2017.03.002}{10.1016/j.jpainsymman.2017.03.002}}

\bibitem[\citeproctext]{ref-verhulst1838}
\CSLLeftMargin{22. }%
\CSLRightInline{Verhulst PF. Notice sur la loi que la population suit dans son accroissement. \emph{Correspondance Math{é}matique et Physique}. 1838;10:113-121.}

\bibitem[\citeproctext]{ref-gompertz1825}
\CSLLeftMargin{23. }%
\CSLRightInline{Gompertz B. On the nature of the function expressive of the law of human mortality, and on a new mode of determining the value of life contingencies. \emph{Philosophical Transactions of the Royal Society of London}. 1825;115:513-583.}

\bibitem[\citeproctext]{ref-winsor1932gompertz}
\CSLLeftMargin{24. }%
\CSLRightInline{Winsor CP. The {G}ompertz curve as a growth curve. \emph{Proceedings of the National Academy of Sciences USA}. 1932;18(1):1-8.}

\bibitem[\citeproctext]{ref-muggeo2003}
\CSLLeftMargin{25. }%
\CSLRightInline{Muggeo VMR. Estimating regression models with unknown break-points. \emph{Statistics in Medicine}. 2003;22(19):3055-3071. doi:\href{https://doi.org/10.1002/sim.1545}{10.1002/sim.1545}}

\end{CSLReferences}

\vfill

\begin{center}\rule{0.5\linewidth}{0.5pt}\end{center}

\emph{Rendered on 2026-07-30 at 11:57 PDT.}\\
\emph{Source: \texttt{\textasciitilde{}/prj/res/36-pmsimstats-ng/pmsimstats-ng/analysis/report/07-gompertz-evaluation/report.Rmd}}

\end{document}
